\documentclass[11pt]{article}
\usepackage[margin=1in]{geometry}
\usepackage{amsmath, amssymb, amsthm}
\usepackage{bm}
\usepackage{hyperref}
\usepackage{enumitem}
\usepackage{booktabs}
\usepackage{listings}
\usepackage[T1]{fontenc}
\usepackage[scaled]{beramono}
\usepackage{courier}
\usepackage{graphicx}

\lstdefinestyle{code}{
  basicstyle=\ttfamily\small,
  breaklines=true,
  frame=single,
  columns=fullflexible
}

\title{%
ACE--SCN / CSC / PhaseMirror-HQ Integration\\[4pt]
\large Technical Report and Defensive Publication
}

\author{Citizen Gardens \\ The Foundation of Multiplicity}
\date{April 25, 2026}

\begin{document}
\maketitle

\begin{abstract}
This document consolidates and formalizes recent developments around the Arithmetic Control Engine (ACE), the Structured Control Network (SCN), the Contraction Spectral Certificate (CSC), and their proposed unification with a PhaseMirror-HQ memory-kernel architecture. It is written as both a technical report and a prior-art defensive publication. The report covers: (i) executive-level motivation and system architecture; (ii) a comprehensive mathematical overview of ACE--SCN, including feasibility maps and certified bounds; (iii) the Track A / Track B split and the canonical 5{,}087-constraint zero-knowledge circuit layout; (iv) the proposal to treat an FZS-MK memory kernel and Zeno projection as the sole semantic authority for drift and protection metrics; (v) an ADR-centered governance plan with SWOT and RAID analysis; and (vi) implementation sketches and validation pathways sufficient to serve as public technical disclosure.%
\end{abstract}
\newpage
\tableofcontents
\newpage
\section{Executive Summary}

The Arithmetic Control Engine (ACE) is a reproducible system for amortized, certified spectral shaping on families of arithmetic operators, in particular self-adjoint Hecke operators on cusp forms.\footnote{Formal specification in ACE--SCN document. [file:1]} Given a Hermitian operator matrix \(A\) and a target gap value \(\hat g\), ACE produces a Hermitian perturbation \(\Delta\) that satisfies hard feasibility certificates and, optionally, preserves arithmetic structure via explicit fidelity modes. The Structured Control Network (SCN) proposes such perturbations in an amortized fashion, while a deterministic feasibility map enforces admissible constraints independently of training.\footnote{Definitions of the feasibility map and fidelity modes are given in Sections 4--6 of the ACE--SCN specification. [file:1]}

The implementation roadmap partitions ACE into two tracks. Track A (Python) computes semantic quantities including drift metrics, windowed contraction (WAC), and certification flags, while Track B (Rust/Circom) implements a fixed-topology zero-knowledge circuit that verifies algebraic invariants, Poseidon2 hash commitments, and an exact canonical constraint budget of 5{,}087 R1CS constraints.\footnote{ACE Roadmap and constraint layout documents. [file:2][file:3]} The file \texttt{csc.py} in Track A is designated as the canonical constraint-budget authority, and the Circom circuit enforces both the Poseidon2 topology \((t=9, r=8)\) and the total constraint count as non-negotiable.%

Within this design thread, the following additional developments have been articulated and are captured here as defensive publication:

\begin{itemize}[leftmargin=*]
  \item Clarification of SCN as the learned amortized controller and CSC (``Contraction Spectral Certificate'') as the topology-and-certificate layer binding SCN outputs, governance witnesses, and cryptographic proofs.
  \item A proposal to unify all drift and protection semantics under a PhaseMirror-HQ memory-kernel layer (FZS-MK with Zeno projection), such that ACE, SCN, and CSC treat this kernel as the sole authority for drift/protection metrics.
  \item An ADR-based governance and migration plan, including SWOT and RAID analyses, to deprecate duplicate ACE-side drift logic while preserving reproducibility and circuit topology constraints.
  \item A telemetry contract for exposing kernel-certified scalars (e.g., normalized drift \(X_n\), maximum WAC statistic, a protection index, and a validity bit) into ACE's certificate and zero-knowledge pipelines, without recomputing semantic logic in the consuming stack.
\end{itemize}

The final architecture proposed is four-layered:
\begin{enumerate}[leftmargin=*]
  \item Arithmetic operator layer (cusp-form spaces, Hecke operators, Hermitianization).
  \item Kernel authority layer (FZS-MK memory kernel and Zeno projection).
  \item Control and certificate layer (SCN, feasibility maps, governance logic, CSC).
  \item Cryptographic verification layer (Circom/Rust circuit, Poseidon2 commitments, Groth16 proofs).
\end{enumerate}

This configuration preserves ACE's original mathematical guarantees while concentrating semantic authority for drift and protection metrics in a single, well-defined kernel.

\section{Mathematical Overview of ACE--SCN}

\subsection{Arithmetic Objects and Hermitianization}

\paragraph{Cusp-form space and safe Hecke family.} Let
\[
V = S_k(\Gamma_0(N))
\]
be the complex vector space of weight-\(k\) holomorphic cusp forms of level \(\Gamma_0(N)\).\footnote{Section 1 of the ACE--SCN specification. [file:1]} For primes \(p\nmid N\), the Hecke operators
\[
T_p : V \to V
\]
are self-adjoint with respect to the Petersson inner product and form the ``safe'' part of the Hecke family; non-self-adjoint operators such as \(U_p\) for \(p\mid N\) are excluded from the base specification.

\paragraph{Computational basis and Gram matrix.} Choose a non-diagonal computable basis \(B\) for \(V\) (for example, a modular-symbols-derived basis), producing a matrix representation
\[
M_p
\]
for \(T_p\) in that basis.\footnote{Section 2 of the ACE--SCN specification. [file:1]} Let \(G\succ 0\) denote the Gram matrix of the relevant inner product in basis \(B\). Self-adjointness in this basis is equivalent to the condition
\[
M_p G = G M_p^*.
\]
A diagnostic self-adjointness error is defined by
\[
\operatorname{selfAdjErr}(M_p,G)
  := \frac{\|M_p G - G M_p^*\|_F}{\|M_p\|_F \,\|G\|_F}.
\]

\paragraph{Stable Hermitianization.} Compute a Cholesky factorization \(G = LL^*\) and define
\[
A_p := L^{-1} M_p L.
\]
If \(\operatorname{selfAdjErr}(A_p)\) is sufficiently small, then \(A_p\) is Hermitian up to numerical tolerance.\footnote{Hermitianization and diagnostics described in Section 2. [file:1]} Optionally, apply a symmetrization
\[
A_p \leftarrow \tfrac12 (A_p + A_p^*)
\]
and record the residual \(\|A_p - A_p^*\|_F\) as a diagnostic.

\subsection{Gap Functionals and Perturbation Classes}

\paragraph{Eigenvalues and gap functionals.} For a Hermitian matrix \(A\in\mathbb{C}^{d\times d}\), denote its ordered real eigenvalues by
\[
\lambda_1(A) \le \cdots \le \lambda_d(A).
\]
The indexed gap is defined as
\[
g_i(A) := \lambda_{i+1}(A) - \lambda_i(A),
\]
while a smooth soft-min gap functional can be defined as
\[
g_\tau(A) := -\tau \log\!\left(\sum_{i=1}^{d-1} \exp\!\left(-\frac{\lambda_{i+1}(A) - \lambda_i(A)}{\tau}\right)\right), \quad \tau > 0.
\]
These functionals provide smooth objectives for spectral-gap control.\footnote{Gap definitions and smooth variants are discussed in Section 3. [file:1]}

\paragraph{Hecke-span subspace and norms.} Fix primes \(p_1,\dots,p_r\) with \(p_j\nmid N\). Define the Frobenius subspace
\[
H_r := \operatorname{span}\{A_{p_1},\dots,A_{p_r}\} \subset \mathbb{C}^{d\times d}.
\]
The Frobenius norm is
\[
\|X\|_F := \sqrt{\operatorname{tr}(X^*X)},
\]
and the induced distance to \(H_r\) is
\[
\operatorname{dist}_F(X,H_r) := \min_{H\in H_r}\|X - H\|_F.
\]

\paragraph{Fidelity modes and constraint sets.} For \(\varepsilon>0\) and \(\eta\ge 0\), ACE defines three constraint sets:
\begin{align*}
\mathcal{C}_\varepsilon
  &:= \{\Delta : \Delta = \Delta^*,\, \|\Delta\|_F \le \varepsilon\}, \\
\mathcal{C}_\varepsilon^{H_r}
  &:= \{\Delta : \Delta\in H_r,\, \Delta = \Delta^*,\, \|\Delta\|_F \le \varepsilon\}, \\
\mathcal{C}^{\mathrm{near}}_{\varepsilon,\eta}
  &:= \{\Delta : \Delta = \Delta^*,\, \|\Delta\|_F \le \varepsilon,\, \operatorname{dist}_F(\Delta,H_r)\le \eta\}.
\end{align*}
These sets correspond respectively to Mode~1 (unstructured), Mode~2 (Hecke-span), and Mode~3 (near-arithmetic) perturbations.\footnote{Fidelity modes appear in Section 4 of the ACE--SCN specification. [file:1]}

\subsection{Projection and Feasibility Map}

\paragraph{Frobenius projection onto \(H_r\).} Define
\[
B = [\operatorname{vec}(A_{p_1}) \; \cdots \; \operatorname{vec}(A_{p_r})] \in \mathbb{R}^{d^2\times r}.
\]
Then the Frobenius projection of a matrix \(X\) onto \(H_r\) is
\[
P_{H_r}(X) = \operatorname{unvec}\bigl(B(B^\top B)^{-1} B^\top \operatorname{vec}(X)\bigr),
\]
with the practical recommendation to precompute a numerically stable orthonormal basis for \(H_r\) via QR factorization on flattened matrices.\footnote{Implementation notes in Section 5. [file:1]}

\paragraph{Certificate-first feasibility map.} Let \(\Delta_0\) be a raw Hermitian proposal (for example, an SCN output after symmetrization). ACE applies a deterministic map \(F\) to \(\Delta_0\), depending on the chosen mode.

\begin{itemize}[leftmargin=*]
  \item Common Hermitianization:
  \[
  \Delta_0 \leftarrow \tfrac12(\Delta_0 + \Delta_0^*).
  \]
  \item Mode 1 (unstructured): scale to the ball of radius \(\varepsilon\),
  \[
  F_1(\Delta_0) := \varepsilon \cdot \frac{\Delta_0}{\max(\|\Delta_0\|_F,\, \varepsilon)}.
  \]
  \item Mode 2 (Hecke-span): project then scale,
  \begin{align*}
    H &= P_{H_r}(\Delta_0), \\
    F_2(\Delta_0) &:= \varepsilon \cdot \frac{H}{\max(\|H\|_F,\, \varepsilon)}.
  \end{align*}
  \item Mode 3 (near-arithmetic): clip residual distance and then scale,
  \begin{align*}
    H &= P_{H_r}(\Delta_0), \\
    R &= \Delta_0 - H, \\
    R' &:= \min\!\left(1,\, \frac{\eta}{\|R\|_F}\right) R, \\
    \Delta_1 &:= H + R', \\
    F_3(\Delta_0) &:= \varepsilon \cdot \frac{\Delta_1}{\max(\|\Delta_1\|_F,\, \varepsilon)}.
  \end{align*}
\end{itemize}

The resulting map \(F_m(\Delta_0)\) (for \(m\in\{1,2,3\}\)) is Hermitian and guaranteed to lie in the corresponding constraint set by construction.\footnote{Feasibility map construction in Section 6. [file:1]}

\subsection{Certificates and Key Lemmas}

\paragraph{Weyl via Frobenius.} For Hermitian \(A\) and Hermitian \(\Delta\in\mathcal{C}_\varepsilon\), Weyl's theorem yields the Frobenius-based bound
\[
|\lambda_i(A+\Delta) - \lambda_i(A)| \le \|\Delta\|_2 \le \|\Delta\|_F \le \varepsilon.
\]
This provides a certified upper bound on eigenvalue movement for all fidelity modes.

\paragraph{Near-span implies near-commuting.} If \(\Delta\) satisfies \(\operatorname{dist}_F(\Delta,H_r)\le \eta\), then for each generator \(A_{p_i}\),
\[
\|[A_{p_i},\Delta]\|_F \le 2\|A_{p_i}\|_2\,\eta.
\]
This inequality turns commutator norms into explicit diagnostics of arithmetic fidelity: small Frobenius distance to the Hecke span implies small commutators.%

\subsection{SCN Amortized Controller}

\paragraph{Dimension-agnostic parametrization.} The SCN learns an amortized mapping from instance features to either a raw matrix proposal or a low-dimensional parameter vector \((\alpha,\beta)\). A recommended parametrization is
\[
\Delta_0(\alpha,\beta)
  = \sum_{j=1}^r \alpha_j A_{p_j}
    + \sum_{\ell=1}^s \beta_\ell u_\ell u_\ell^\top,
\]
where \(u_\ell \in \mathbb{R}^d\) are deterministic pseudorandom unit vectors generated from a fixed seed. This makes the network output dimension independent of \(d\), the ambient matrix size.

\paragraph{Features and amortized objective.} Let \(\phi(A,\hat g)\) be a feature vector including target gap \(\hat g\), eigenvalue summaries, selected gap values, and quantiles. For a fixed constraint set \(\mathcal{C}\), the per-instance constrained problem is
\[
\Delta^*(A,\hat g)
  \in \arg\min_{\Delta\in\mathcal{C}}\left(g_\tau(A+\Delta) - \hat g\right)^2.
\]
The amortized training objective for SCN parameters \(\theta\) is
\[
\min_\theta \mathbb{E}_{(A,\hat g)}\left[
  \left(
    g_\tau\bigl(A + F(\Delta_0(F_\theta(\phi(A,\hat g))))\bigr) - \hat g
  \right)^2
\right],
\]
where \(F_\theta\) is the network mapping features to parameter vectors, \(\Delta_0(\cdot)\) constructs raw perturbations, and \(F(\cdot)\) is the deterministic feasibility map.

\section{Governance Metrics and Zero-Knowledge Layout}

\subsection{Track A Governance Metrics}

The ACE roadmap introduces a governance stack that computes various drift and retention metrics in Track A.\footnote{Track A description in ACE Roadmap. [file:2]} Given a proxy distribution \(\mu_{\mathrm{proxy}}\) and a current empirical mean \(\mu_n\), define:

\begin{itemize}[leftmargin=*]
  \item Proxy \(L_1\) norm:
  \[
  \|\mu_{\mathrm{proxy}}\|_1 = \sum_i |\mu_{\mathrm{proxy},i}|.
  \]
  \item Support mask and off-support drift:
  \[
  \bar m_n = \sum_i |\mu_n(i) - \mu_{\mathrm{proxy}}(i)|\,(1 - \mathrm{support\_mask}_i).
  \]
  \item Normalized drift metric:
  \[
  X_n = \frac{\bar m_n}{\|\mu_{\mathrm{proxy}}\|_1}.
  \]
  \item Retention rate:
  \[
  R_t = \max(0, 1 - \varepsilon - X_n).
  \]
  \item Telescoping WAC window metric:
  \[
  W_t = \frac{\bar m^{(t+M)}}{\max(\bar m^{(t)},10^{-12})}.
  \]
\end{itemize}

The flag \(\mathrm{is\_valid}\in\{0,1\}\) is set based on whether \(\max_t W_t\) stays below a prescribed threshold (typically less than 1), serving as a certification gate for windowed contraction.

\subsection{Track B Circuit and Constraint Budget}

The Track B circuit, implemented in Circom and orchestrated by Rust, is constrained by a canonical 5{,}087-constraint budget that is locked both in \texttt{csc.py} and in a constraint-layout specification.\footnote{Constraint layout document. [file:3]} The optimized base circuit uses 133 constraints, decomposed into:
\begin{itemize}[leftmargin=*]
  \item 64 trajectory consistency constraints:
  \[
  \mathrm{current\_mu}[i] \cdot \mathrm{step\_n} \equiv \mathrm{h\_hat}[i] \pmod p.
  \]
  \item 64 support-mask booleanity constraints:
  \[
  \mathrm{support\_mask}[i](\mathrm{support\_mask}[i]-1) \equiv 0 \pmod p.
  \]
  \item 1 drift quotient check:
  \[
  X_n \cdot \mathrm{proxy\_l1\_norm} \equiv \bar m_n\cdot \mathrm{SCALE} \pmod p.
  \]
  \item 1 clamping consistency:
  \[
  R_t \cdot (r_{\mathrm{raw}} - R_t) \equiv 0 \pmod p,
  \quad r_{\mathrm{raw}} = \mathrm{SCALE}-\varepsilon-X_n.
  \]
  \item 1 validity-flag booleanity:
  \[
  \mathrm{is\_valid}(\mathrm{is\_valid}-1) \equiv 0 \pmod p.
  \]
  \item 2 topology guards for WAC mode and \(N_0\) consistency.
\end{itemize}

Poseidon2 hash gadgets (sponge and single-permutation) are then added to reach a projected 4{,}804 constraints after hashing all witness fields and binding commitments, leaving 283 constraints in reserve for guard-rails and audit hooks.\footnote{Poseidon2 integration roadmap and budget in constraint layout. [file:3]}

\subsection{Commitment Strategy}

Track B uses two Poseidon2 layers:
\begin{enumerate}[leftmargin=*]
  \item A sponge hash over all relevant witness fields, producing a commitment \(h_{\mathrm{commitment}}\).
  \item A second Poseidon2 permutation binding the hash and selected outputs (e.g., \(X_n\), \(R_t\), \(\mathrm{is\_valid}\), WAC product, and a retry nonce) into a final CAS commitment.
\end{enumerate}
In Circom-style pseudocode:
\begin{lstlisting}[style=code,language=C]
component poseidon_h = Poseidon2Sponge(205, 9, 8);
for (var i = 0; i < 205; i++) {
    poseidon_h.in[i] <== witness_field[i];
}
signal h_commitment;
h_commitment <== poseidon_h.out;

component poseidon_gamma = Poseidon2(5, 9, 8);
poseidon_gamma.in[0] <== h_commitment;
poseidon_gamma.in[1] <== X_n;
poseidon_gamma.in[2] <== R_t;
poseidon_gamma.in[3] <== max_wac_product;
poseidon_gamma.in[4] <== retry_nonce;
cas_commitment <== poseidon_gamma.out;
\end{lstlisting}
This structure is tightly coupled to the 5{,}087-constraint budget and enforces both determinism and collision resistance.\footnote{Detailed in the constraint layout and roadmap. [file:2][file:3]}

\section{PhaseMirror-HQ Unification and Kernel Authority}

\subsection{Motivation for Single Semantic Authority}

In the current ACE design, Track A defines the semantics of drift and protection metrics (e.g.\ \(X_n\), WAC, and \(\mathrm{is\_valid}\)), while Track B only verifies algebraic consistency and commitments. This leaves open the possibility of duplicating or partially re-implementing drift logic in other systems interacting with ACE.

The proposed unification introduces a dedicated \emph{kernel authority layer}, implemented by PhaseMirror-HQ, whose FZS-MK memory kernel and Zeno projection define all drift and protection metrics consumed by ACE, SCN, and CSC. Under this model:
\begin{itemize}[leftmargin=*]
  \item Semantic definitions of drift and protection reside entirely in the kernel layer.
  \item ACE and SCN treat these metrics as read-only telemetry.
  \item CSC verifies that these telemetry values are properly bound into commitments and respect topological and budget constraints, but does not define their semantics.
\end{itemize}

\subsection{Kernel Telemetry Contract}

A minimal telemetry record emitted by PhaseMirror could be
\[
\Theta_{\mathrm{kernel}} = \left(
  X_n^{\mathrm{kernel}},
  W_{\max}^{\mathrm{kernel}},
  \zeta_{\mathrm{protect}},
  \mathrm{is\_valid}_{\mathrm{kernel}}
\right),
\]
augmented with a version field for schema evolution. ACE would then consume this record as control telemetry and include its fields in certificate and CAS commitments.

A Python sketch of this interface is:

\begin{lstlisting}[style=code,language=Python]
from dataclasses import dataclass

@dataclass(frozen=True)
class KernelTelemetry:
    xn_kernel: float
    wt_max_kernel: float
    protection_zeta: float
    is_valid_kernel: bool
    telemetry_version: int = 1

def load_kernel_telemetry(step_state) -> KernelTelemetry:
    # Delegates semantic authority to PhaseMirror-HQ/kernel/core/zeno
    return zeno.compute_kernel_telemetry(step_state)
\end{lstlisting}

An ACE-side wrapper can optionally maintain a \emph{parity mode} during migration:

\begin{lstlisting}[style=code,language=Python]
def get_drift_metrics(step_state, parity_mode: bool = False):
    kt = load_kernel_telemetry(step_state)
    if parity_mode:
        legacy_xn = compute_legacy_xn(step_state)
        assert abs(legacy_xn - kt.xn_kernel) < 1e-9
    return kt
\end{lstlisting}

\subsection{SCN Conditioning on Kernel Telemetry}

With kernel authority established, the SCN feature map is extended from \(\phi(A,\hat g)\) to
\[
\phi(A,\hat g, \Theta_{\mathrm{kernel}}),
\]
allowing the controller to react to kernel-defined drift and protection states without recomputing or redefining them. The feasibility map and spectral-control objective remain unchanged; only the policy distribution learned by SCN is conditioned on the additional telemetry.

\section{ADR, SWOT, and RAID Structure}

\subsection{ADR Package}

An ADR (Architecture Decision Record) package for this unification can be organized as:
\begin{itemize}[leftmargin=*]
  \item \textbf{ADR--001}: FZS-MK/Zeno as Single Drift and Protection Authority.
  \item \textbf{ADR--002}: ACE--PhaseMirror Kernel Telemetry Contract.
  \item \textbf{ADR--003}: SCN/CSC Conditioning on Kernel Metrics Only.
\end{itemize}

Each ADR should include context, decision, rationale, alternatives, consequences, SWOT, and a RAID log. A skeletal structure (in prose form) is:

\subsubsection*{ADR--001: FZS-MK/Zeno as Single Drift and Protection Authority}

\begin{itemize}[leftmargin=*]
  \item \textbf{Context}: Drift and protection metrics currently defined in ACE Track A; PhaseMirror kernel provides a potentially richer and more general non-Markovian framework.
  \item \textbf{Decision}: Declare PhaseMirror's FZS-MK/Zeno layer the sole semantic authority for drift and protection metrics used by ACE, SCN, and CSC.
  \item \textbf{Consequences}: ACE may consume and certify these metrics but must not define independent semantic formulas after migration completes.
\end{itemize}

Similar detail can be elaborated for ADR--002 and ADR--003, focusing on schema-level contracts and SCN feature constraints, respectively.

\subsection{SWOT Analysis}

\paragraph{Strengths.}
\begin{itemize}[leftmargin=*]
  \item ACE already separates learned proposal generation from deterministic feasibility enforcement, making authority unification structurally compatible.\footnote{SCN and feasibility map design. [file:1]}
  \item The Track A / Track B split ensures that semantic computation resides outside the circuit, simplifying kernel integration.\footnote{Roadmap and constraint layout. [file:2][file:3]}
  \item A fixed 5{,}087-constraint budget with Poseidon2 topology lock yields strong reproducibility and auditability.\footnote{Constraint layout. [file:3]}
\end{itemize}

\paragraph{Weaknesses.}
\begin{itemize}[leftmargin=*]
  \item Current documentation does not yet specify the exact ACE--PhaseMirror API boundary.
  \item Moving semantics to PhaseMirror increases coupling to kernel release cycles and failure modes.
  \item The circuit relies on Track A for semantic correctness, so kernel bugs are not directly ruled out by the R1CS constraints.\footnote{Constraint responsibilities. [file:3]}
\end{itemize}

\paragraph{Opportunities.}
\begin{itemize}[leftmargin=*]
  \item Prevents semantic drift across training, runtime, and proof-generation pipelines.
  \item Telemetry contract enables both open-core and OEM deployment paths without altering high-level mathematics.\footnote{Deployment discussion in roadmap. [file:2]}
  \item Reserved circuit headroom (283 constraints) allows introducing future guard rails or audit hooks around kernel telemetry if needed.\footnote{Constraint budget headroom. [file:3]}
\end{itemize}

\paragraph{Threats.}
\begin{itemize}[leftmargin=*]
  \item Divergent telemetry schemas between ACE and PhaseMirror could break determinism or governance clarity.
  \item Overly complex telemetry could strain Poseidon2 input layout or circuit budget.
  \item Partial or inconsistent deprecation of ACE-side drift logic could undermine the single-authority guarantee.
\end{itemize}

\subsection{RAID Log (Risks, Assumptions, Issues, Dependencies)}

\paragraph{Risks.}
\begin{itemize}[leftmargin=*]
  \item \textbf{R1}: Numerical divergence between kernel and legacy ACE drift metrics during migration.
  \item \textbf{R2}: Telemetry schema churn compromising certificate determinism.
  \item \textbf{R3}: Telemetry additions inadvertently increasing constraint count beyond budget.
\end{itemize}

\paragraph{Assumptions.}
\begin{itemize}[leftmargin=*]
  \item \textbf{A1}: PhaseMirror kernel can produce stable, scalar (or low-dimensional) telemetry suitable for ACE.
  \item \textbf{A2}: Circuit headroom is sufficient for any minor witness-binding refinements.\footnote{Headroom explicitly reserved for future guard-rails. [file:3]}
\end{itemize}

\paragraph{Issues.}
\begin{itemize}[leftmargin=*]
  \item \textbf{I1}: Formalizing a versioned \texttt{KernelTelemetry} schema.
  \item \textbf{I2}: Establishing a deprecation policy and test coverage for legacy ACE drift/WAC logic.
\end{itemize}

\paragraph{Dependencies.}
\begin{itemize}[leftmargin=*]
  \item \textbf{D1}: PhaseMirror-HQ kernel API release.
  \item \textbf{D2}: ACE certificate schema updates and Track B witness adapter changes.
\end{itemize}

\section{Code Snippets and Implementation Sketches}

\subsection{Python: Kernel Telemetry and Certificate Payload}

\begin{lstlisting}[style=code,language=Python]
from dataclasses import dataclass

@dataclass(frozen=True)
class KernelTelemetry:
    xn_kernel: float
    wt_max_kernel: float
    protection_zeta: float
    is_valid_kernel: bool
    telemetry_version: int = 1

def load_kernel_telemetry(step_state) -> KernelTelemetry:
    # Delegates semantic authority to PhaseMirror-HQ/kernel/core/zeno
    return zeno.compute_kernel_telemetry(step_state)

def get_drift_metrics(step_state, parity_mode: bool = False):
    kt = load_kernel_telemetry(step_state)
    if parity_mode:
        legacy_xn = compute_legacy_xn(step_state)
        assert abs(legacy_xn - kt.xn_kernel) < 1e-9
    return kt

def certificate_payload(theta_base, telemetry, outputs):
    return {
        "theta": serialize_theta(theta_base),
        "xn_kernel": telemetry.xn_kernel,
        "wt_max_kernel": telemetry.wt_max_kernel,
        "protection_zeta": telemetry.protection_zeta,
        "is_valid_kernel": int(telemetry.is_valid_kernel),
        "telemetry_version": telemetry.telemetry_version,
        "outputs": outputs,
    }
\end{lstlisting}

\subsection{Circom Witness Binding Sketch}

Adapting the existing CAS commitment pattern, a telemetry-aware binding might be:

\begin{lstlisting}[style=code]
component poseidon_gamma = Poseidon2(5, 9, 8);
poseidon_gamma.in[0] <== h_commitment;
poseidon_gamma.in[1] <== xn_kernel;
poseidon_gamma.in[2] <== retention_rate;
poseidon_gamma.in[3] <== max_wac_product;
poseidon_gamma.in[4] <== retry_nonce;
cas_commitment <== poseidon_gamma.out;
\end{lstlisting}

Here, \texttt{xn\_kernel} replaces or augments the previous \(X_n\) witness field while preserving circuit budget and topology constraints.

\section{Validation Path and Repository Changes}

\subsection{Four-Phase Validation Plan}

A fast yet robust validation path is:

\begin{enumerate}[leftmargin=*]
  \item \textbf{Authority Declaration}: Accept ADR--001 and update documentation to state that PhaseMirror's kernel layer is the semantic authority for drift and protection metrics.
  \item \textbf{Parity Phase}: Run existing ACE drift/WAC logic in parallel with kernel telemetry on benchmark datasets and assert numerical agreement within strict tolerances.
  \item \textbf{Contract Phase}: Integrate \texttt{KernelTelemetry} fields into ACE certificates and Poseidon2 witness commitments, keeping SCN behavior unchanged initially.
  \item \textbf{Control Phase}: Retrain or fine-tune SCN to condition on kernel telemetry, then compare to pre-integration baselines in terms of objective values, fidelity metrics, proof determinism, and runtime.
\end{enumerate}

\subsection{Recommended Repository Edits}

Within the ACE repository structure outlined in the roadmap,\footnote{ACE Core scaffold and modules. [file:2]} recommended edits include:

\begin{itemize}[leftmargin=*]
  \item \textbf{\texttt{src/ace/governor.py}}: Replace internal ownership of \(X_n\) with calls into \texttt{get\_drift\_metrics}.
  \item \textbf{\texttt{src/ace/monitor.py}}: Route WAC-related decisions via kernel telemetry.
  \item \textbf{\texttt{src/ace/certify.py}}: Extend the certificate schema with telemetry fields and versioning.
  \item \textbf{\texttt{src/ace/cas\_registry.py}}: Ensure telemetry fields are hashed consistently into CAS commitments.
  \item \textbf{\texttt{src/ace/csc.py}} and Circom circuit: Maintain the 5{,}087-constraint lock and adjust witness binding if necessary within the existing budget.
  \item \textbf{\texttt{tests/}}: Add parity tests (kernel vs legacy drift), schema-version tests, and checks that forbid reintroduction of duplicate drift logic.
  \item \textbf{\texttt{adr/}}: Add ADR--001, ADR--002, and ADR--003, each with its own embedded SWOT and RAID sections.
\end{itemize}

\section{Defensive Publication Notes}

This report is intended as a defensive disclosure of the integrated ACE--SCN--CSC--PhaseMirror architecture. It documents:
\begin{itemize}[leftmargin=*]
  \item The mathematical specification of ACE--SCN, including fidelity modes, feasibility maps, and certified commutator bounds.
  \item The Track A / Track B split and the fixed 5{,}087-constraint Circom topology with Poseidon2 commitments.
  \item The conceptual role of a Contraction Spectral Certificate as a certificate-and-topology layer binding controller outputs to cryptographic proofs.
  \item The unification with a memory-kernel and projection-based authority layer for drift and protection metrics.
  \item Concrete code and interface sketches for telemetry integration, ADR governance structures, SWOT and RAID analyses, and validation pathways.
\end{itemize}

For maximum defensive value, it is recommended to:
\begin{itemize}[leftmargin=*]
  \item Publish this document (or an equivalent) in a publicly timestamped venue or repository.
  \item Include immutable hashes or content-addressable identifiers for this text and associated ADRs.
  \item Maintain accompanying machine-readable schemas (e.g., JSON schemas for \texttt{KernelTelemetry} and \texttt{ACECertificate}) in a public code repository.
\end{itemize}

\section{Conclusion}

The developments synthesized here refine rather than replace the ACE program. ACE remains a certified spectral-control engine on arithmetic operators; SCN remains the amortized proposal mechanism; CSC remains the topology-and-certificate authority; and a PhaseMirror FZS-MK/Zeno layer becomes the single semantic authority for drift and protection metrics. This layered arrangement is mathematically coherent, operationally auditable, and appropriate for a prior-art defensive publication.

\appendix
\section{Mathematical Appendix: Proofs and Operator Norm Bounds}

This appendix collects explicit proofs and operator-norm bounds underlying the ACE--SCN specification, with emphasis on feasibility-map invariants, spectral perturbation bounds, and commutator estimates used as arithmetic-fidelity diagnostics.\footnote{See Sections 4--9 of the ACE--SCN formal specification for the informal statements. [file:1]} Throughout, \( \|\cdot\|_2 \) denotes the operator norm induced by the Euclidean norm on \( \mathbb{C}^d \), and \( \|\cdot\|_F \) denotes the Frobenius norm.

\subsection{Weyl-Type Bounds via Frobenius Norm}

ACE uses Weyl's inequality together with the relation \( \|\cdot\|_2 \le \|\cdot\|_F \) to certify that any perturbation in \( \mathcal{C}_\varepsilon \) moves eigenvalues by at most \( \varepsilon \) in absolute value.\footnote{The set \( \mathcal{C}_\varepsilon \) is defined in Section 4.3 of the ACE--SCN document. [file:1]}

\begin{lemma}[Weyl via Frobenius]\label{lem:weyl-frob}
Let \( A \in \mathbb{C}^{d\times d} \) be Hermitian and let \( \Delta \in \mathbb{C}^{d\times d} \) be Hermitian with \( \|\Delta\|_F \le \varepsilon \). Denote the ordered eigenvalues of \( A \) and \( A+\Delta \) by
\[
\lambda_1(A) \le \cdots \le \lambda_d(A), \qquad
\lambda_1(A+\Delta) \le \cdots \le \lambda_d(A+\Delta).
\]
Then for each \( i \in \{1,\dots,d\} \),
\[
|\lambda_i(A+\Delta) - \lambda_i(A)| \;\le\; \|\Delta\|_2 \;\le\; \|\Delta\|_F \;\le\; \varepsilon.
\]
\end{lemma}

\begin{proof}
The first inequality is Weyl's eigenvalue perturbation bound for Hermitian matrices:
\[
|\lambda_i(A+\Delta) - \lambda_i(A)| \le \|\Delta\|_2, \quad i=1,\dots,d.
\]
This is a standard result; see, e.g., classical matrix analysis texts. The second inequality
\[
\|\Delta\|_2 \le \|\Delta\|_F
\]
follows because the Frobenius norm is the \( \ell_2 \)-norm of the singular values, and the operator norm is the \( \ell_\infty \)-norm of the same singular values, so \( \|\cdot\|_2 \le \|\cdot\|_F \) always holds. The final inequality is by the assumption \( \|\Delta\|_F \le \varepsilon \), completing the chain.
\end{proof}

As an immediate corollary, any perturbation produced by the feasibility map in Mode~1 (which outputs an element of \( \mathcal{C}_\varepsilon \)) moves every eigenvalue of the target operator \( A \) by at most \( \varepsilon \).

\subsection{Feasibility-Map Invariants in All Modes}

We next give a unified proof that the feasibility maps \( F_1, F_2, F_3 \) used in ACE preserve Hermitian symmetry and enforce the appropriate mode constraints.\footnote{See Section 6 of the ACE--SCN specification for the construction of \( F_m \). [file:1]}

\begin{proposition}[Feasibility-map invariants]\label{prop:feasibility}
Let \( \Delta_0 \in \mathbb{C}^{d\times d} \) be arbitrary, and let
\[
\Delta_0' := \tfrac12(\Delta_0 + \Delta_0^*)
\]
be its Hermitian part. Let \( F_1, F_2, F_3 \) be defined as in the ACE specification, using \( \Delta_0' \) as input:
\begin{align*}
F_1(\Delta_0') &:= \varepsilon \cdot \frac{\Delta_0'}{\max(\|\Delta_0'\|_F,\,\varepsilon)}, \\
H(\Delta_0') &:= P_{H_r}(\Delta_0'), \\
F_2(\Delta_0') &:= \varepsilon \cdot \frac{H(\Delta_0')}{\max(\|H(\Delta_0')\|_F,\,\varepsilon)}, \\
R(\Delta_0') &:= \Delta_0' - H(\Delta_0'), \\
R'(\Delta_0') &:= \min\!\left(1,\,\frac{\eta}{\|R(\Delta_0')\|_F}\right)R(\Delta_0'), \\
\Delta_1(\Delta_0') &:= H(\Delta_0') + R'(\Delta_0'), \\
F_3(\Delta_0') &:= \varepsilon \cdot \frac{\Delta_1(\Delta_0')}{\max(\|\Delta_1(\Delta_0')\|_F,\,\varepsilon)}.
\end{align*}
Then:
\begin{enumerate}[label=(\alph*), leftmargin=*]
  \item For each \( m\in\{1,2,3\} \), \( F_m(\Delta_0') \) is Hermitian.
  \item \( F_1(\Delta_0') \in \mathcal{C}_\varepsilon \).
  \item \( F_2(\Delta_0') \in \mathcal{C}_\varepsilon^{H_r} \).
  \item \( F_3(\Delta_0') \in \mathcal{C}^{\mathrm{near}}_{\varepsilon,\eta} \).
\end{enumerate}
\end{proposition}

\begin{proof}
(a) Hermitian symmetry is preserved in each case because all operations used in the construction are Hermitian-preserving:

\begin{itemize}
  \item \( \Delta_0' \) is Hermitian by definition.
  \item The Frobenius projection \( P_{H_r} \) onto a linear subspace \( H_r \) preserves Hermitian structure when \( H_r \) is spanned by Hermitian matrices \( A_{p_j} \). That is, if \( \Delta_0' = \Delta_0'^* \), then \( H(\Delta_0') \) is a real linear combination of the \( A_{p_j} \), hence Hermitian.
  \item The residual \( R(\Delta_0') = \Delta_0' - H(\Delta_0') \) is the difference of Hermitian matrices, hence Hermitian.
  \item Scaling by a real scalar preserves Hermitian structure, so \( R'(\Delta_0') \) and the normalized outputs \( F_m(\Delta_0') \) remain Hermitian.
\end{itemize}

(b) For Mode~1, define
\[
\Delta^{(1)} := F_1(\Delta_0') = \varepsilon \cdot \frac{\Delta_0'}{\max(\|\Delta_0'\|_F,\,\varepsilon)}.
\]
By construction, \( \Delta^{(1)} = (\Delta^{(1)})^* \). Its Frobenius norm is
\[
\|\Delta^{(1)}\|_F
  = \varepsilon \cdot \frac{\|\Delta_0'\|_F}{\max(\|\Delta_0'\|_F,\,\varepsilon)}
  \le \varepsilon,
\]
since the fraction is at most 1. Thus \( \Delta^{(1)} \in \mathcal{C}_\varepsilon \).

(c) For Mode~2, define
\[
\Delta^{(2)} := F_2(\Delta_0') = \varepsilon \cdot \frac{H(\Delta_0')}{\max(\|H(\Delta_0')\|_F,\,\varepsilon)}.
\]
We have \( H(\Delta_0') \in H_r \) by definition of the projection. As above, the scaling factor ensures
\[
\|\Delta^{(2)}\|_F \le \varepsilon.
\]
Hermitianity of \( H(\Delta_0') \) was established in part (a). Therefore, \( \Delta^{(2)} \in \mathcal{C}_\varepsilon^{H_r} \).

(d) For Mode~3, write
\[
\Delta^{(3)} := F_3(\Delta_0') = \varepsilon \cdot \frac{\Delta_1(\Delta_0')}{\max(\|\Delta_1(\Delta_0')\|_F,\,\varepsilon)}.
\]
By part (a), \( \Delta_1(\Delta_0') \) is Hermitian. The final scaling ensures
\[
\|\Delta^{(3)}\|_F \le \varepsilon.
\]

It remains to show that \( \operatorname{dist}_F(\Delta^{(3)}, H_r) \le \eta \). Let
\[
\widetilde{\Delta} := \Delta_1(\Delta_0') = H(\Delta_0') + R'(\Delta_0').
\]
We have \( H(\Delta_0')\in H_r \) and \( R'(\Delta_0') \) orthogonal (in the Frobenius geometry) to \( H_r \) because it is a scaled version of the residual from the projection. Thus
\[
\operatorname{dist}_F(\widetilde{\Delta},H_r) = \|R'(\Delta_0')\|_F.
\]
By definition of \( R'(\Delta_0') \),
\[
\|R'(\Delta_0')\|_F
  = \min\!\left(1,\frac{\eta}{\|R(\Delta_0')\|_F}\right)\|R(\Delta_0')\|_F
  \le \eta.
\]
Now note that the final scaling
\[
\Delta^{(3)} = \varepsilon \cdot \widetilde{\Delta}/\max(\|\widetilde{\Delta}\|_F,\,\varepsilon)
\]
shrinks \(\widetilde{\Delta}\) (or leaves it unchanged) by a nonnegative scalar factor no larger than 1. In particular, distances to \(H_r\) (measured in Frobenius norm) can only decrease under such a contraction:
\[
\operatorname{dist}_F(\Delta^{(3)},H_r)
  \le \operatorname{dist}_F(\widetilde{\Delta},H_r)
  = \|R'(\Delta_0')\|_F
  \le \eta.
\]
Therefore, \( \Delta^{(3)} \in \mathcal{C}^{\mathrm{near}}_{\varepsilon,\eta} \), completing the proof.
\end{proof}

\subsection{Near-Span Implies Near-Commuting}

ACE interprets commutator norms \( \|[A_{p_i},\Delta]\|_F \) as diagnostics for arithmetic fidelity in Mode~3.\footnote{See Section 9.2 of the ACE--SCN specification. [file:1]} Here we spell out the standard operator-norm argument that justifies the bound.

\begin{proposition}[Near-span implies near-commuting]\label{prop:near-comm}
Let \( A_{p_i} \in \mathbb{C}^{d\times d} \) be Hermitian operators (e.g.\ the safe Hecke matrices) spanning \(H_r\), and let
\[
\Delta \in \mathcal{C}^{\mathrm{near}}_{\varepsilon,\eta}
\]
so that \( \Delta = \Delta^* \) and \( \operatorname{dist}_F(\Delta,H_r)\le \eta \). Then for each generator \(A_{p_i}\),
\[
\|[A_{p_i},\Delta]\|_F \;\le\; 2\,\|A_{p_i}\|_2\,\eta.
\]
\end{proposition}

\begin{proof}
By definition of the distance, there exists \(H\in H_r\) such that
\[
\|\Delta - H\|_F = \operatorname{dist}_F(\Delta,H_r) \le \eta.
\]
We can decompose
\[
[A_{p_i},\Delta] = [A_{p_i},H] + [A_{p_i},\Delta - H].
\]
Since \(H\in H_r\) and \(H_r\) is spanned by the commuting family of safe Hecke operators \(A_{p_j}\), \(H\) commutes with each \(A_{p_i}\), giving
\[
[A_{p_i},H] = 0.
\]
Therefore,
\[
[A_{p_i},\Delta] = [A_{p_i},\Delta - H] = A_{p_i}(\Delta-H) - (\Delta-H)A_{p_i}.
\]

Using submultiplicativity of the operator and Frobenius norms,
\begin{align*}
\|[A_{p_i},\Delta]\|_F
  &= \|A_{p_i}(\Delta-H) - (\Delta-H)A_{p_i}\|_F \\
  &\le \|A_{p_i}(\Delta-H)\|_F + \|(\Delta-H)A_{p_i}\|_F \\
  &\le \|A_{p_i}\|_2 \|\Delta-H\|_F + \|\Delta-H\|_F \|A_{p_i}\|_2 \\
  &= 2\,\|A_{p_i}\|_2\,\|\Delta - H\|_F \\
  &\le 2\,\|A_{p_i}\|_2\,\eta,
\end{align*}
as claimed.
\end{proof}

This proposition is the precise content behind ACE's statement that Mode~3 commutator norms scale as \(O(\eta)\) and can be interpreted as certified arithmetic-fidelity diagnostics.

\subsection{Monotonicity and Scaling Behaviour in Mode 3}

ACE informally notes that in Mode~3, smaller \(\eta\) produces stricter control of arithmetic fidelity and smaller commutators while Mode~2 corresponds to the endpoint \(\eta = 0\).\footnote{See Section 12 of the ACE--SCN specification for the qualitative description of expected outcomes. [file:1]} We briefly formalize these monotonicity properties.

\begin{lemma}[Monotonicity in \texorpdfstring{\(\eta\)}{eta}]\label{lem:eta-monotone}
Fix \(\varepsilon>0\), and consider two near-arithmetic constraint sets
\[
\mathcal{C}^{\mathrm{near}}_{\varepsilon,\eta_1}
  \subseteq
\mathcal{C}^{\mathrm{near}}_{\varepsilon,\eta_2}
\quad\text{with}\quad
0 \le \eta_1 \le \eta_2.
\]
Then the set of admissible perturbations under \(\eta_1\) is contained in the admissible set under \(\eta_2\). In particular, any admissible \(\Delta\) for \(\eta_1\) is also admissible for \(\eta_2\), and the corresponding commutator bounds satisfy
\[
\|[A_{p_i},\Delta]\|_F \le 2\|A_{p_i}\|_2\,\eta_1
\quad\Rightarrow\quad
\|[A_{p_i},\Delta]\|_F \le 2\|A_{p_i}\|_2\,\eta_2.
\]
\end{lemma}

\begin{proof}
The inclusion of sets is immediate from the definition:
\[
\mathcal{C}^{\mathrm{near}}_{\varepsilon,\eta_1}
  = \{\Delta: \Delta = \Delta^*,\ \|\Delta\|_F\le\varepsilon,\ \operatorname{dist}_F(\Delta,H_r)\le\eta_1\},
\]
and if \(\eta_1\le \eta_2\), then every \(\Delta\) satisfying \(\operatorname{dist}_F(\Delta,H_r)\le\eta_1\) also satisfies \(\operatorname{dist}_F(\Delta,H_r)\le\eta_2\). The rest follows from Proposition~\ref{prop:near-comm}.
\end{proof}

\begin{corollary}[Mode~2 as endpoint]\label{cor:mode2-endpoint}
Mode~2 corresponds exactly to the case \(\eta=0\) in Mode~3. In this limit, the near-arithmetic constraint set reduces to the Hecke-span constraint set, and the commutator bound becomes
\[
\|[A_{p_i},\Delta]\|_F = 0
\]
for all \(A_{p_i}\), since \(\Delta\in H_r\) commutes with each generator.
\end{corollary}

\begin{proof}
If \(\eta=0\), then \(\operatorname{dist}_F(\Delta,H_r)\le 0\) implies \(\Delta\in H_r\), so
\[
\mathcal{C}^{\mathrm{near}}_{\varepsilon,0} = \mathcal{C}_\varepsilon^{H_r}.
\]
The commutator statement follows as in Proposition~\ref{prop:near-comm} with \(\eta=0\).
\end{proof}

\subsection{Operator Norm Bounds in the Drift Quotient Constraint}

The constraint-layout document encodes the normalized drift metric \(X_n\) and the off-support drift sum \(\bar m_n\) via a single multiplicative constraint
\[
X_n \cdot \mathrm{proxy\_l1\_norm} \equiv \bar m_n \cdot \mathrm{SCALE} \pmod{p}, 
\]
which, in real arithmetic, corresponds to
\[
X_n = \frac{\bar m_n}{\|\mu_{\mathrm{proxy}}\|_1}.
\]\footnote{Normalization constraint as in the ACE constraint layout. [file:3]}

While this is an exact identity rather than an inequality, it interacts with the norm bounds above via the role of \(X_n\) in retention and feasibility constraints. For completeness, we note the following simple bound.

\begin{lemma}[Drift and retention bounds]\label{lem:drift-retention}
Assume the proxy distribution \(\mu_{\mathrm{proxy}}\) is normalized so that \(\|\mu_{\mathrm{proxy}}\|_1 > 0\), and let
\[
\bar m_n = \sum_{i} |\mu_n(i) - \mu_{\mathrm{proxy}}(i)|\,(1-\mathrm{support\_mask}_i)
\]
with \(\mathrm{support\_mask}_i\in\{0,1\}\). Then
\[
0 \le X_n = \frac{\bar m_n}{\|\mu_{\mathrm{proxy}}\|_1}
\]
and the retention rate
\[
R_t = \max(0, 1 - \varepsilon - X_n)
\]
satisfies
\[
0 \le R_t \le 1.
\]
\end{lemma}

\begin{proof}
By definition, each summand in \(\bar m_n\) is nonnegative, hence \(\bar m_n\ge 0\). With \(\|\mu_{\mathrm{proxy}}\|_1>0\), this implies \(X_n\ge 0\). The upper bound \(R_t\le 1\) follows from
\[
R_t = \max(0, 1 - \varepsilon - X_n) \le 1,
\]
since \(\varepsilon>0\) and \(X_n\ge 0\). Nonnegativity is immediate from the max operation.
\end{proof}

This simple bound is one of the ingredients justifying the clamping constraint
\[
R_t(r_{\mathrm{raw}} - R_t) = 0,\quad r_{\mathrm{raw}} = \mathrm{SCALE} - \varepsilon - X_n
\]
in the circuit, which encodes \(R_t = \max(0,r_{\mathrm{raw}})\) in arithmetic form.\footnote{Clamping constraint described in the constraint layout. [file:3]}

\subsection{Summary}

The bounds and proofs in this appendix make explicit the operator-norm control that underpins ACE's spectral guarantees:
\begin{itemize}[leftmargin=*]
  \item Feasibility maps yield perturbations lying in the correct constraint sets while preserving Hermitianity.
  \item Eigenvalue movements are bounded by the Frobenius radius \(\varepsilon\) via Weyl's inequality.
  \item Commutator norms behave as \(O(\eta)\) in Mode~3, with Mode~2 as the commuting endpoint.
  \item Drift and retention metrics used in the Track~A governance stack respect natural range constraints that are encoded succinctly in the Track~B circuit.
\end{itemize}
These explicit derivations complement the main specification and strengthen the correctness claims made in the ACE–SCN and constraint-layout documents.\footnote{See the ACE formal specification and constraint layout for the high-level statements this appendix formalizes. [file:1][file:3]}
\nocite{*}
\bibliographystyle{unsrt}  
\bibliography{references}
\end{document}